% byu-coi-management-plan-example.tex
%
% LaTeX source for the PDF rendering of ../byu-coi-management-plan-example.md
% (vertical-cloud-lab/byu-vcl). The Markdown file is canonical; if it changes,
% update this file to match and recompile:
%
%   pdflatex byu-coi-management-plan-example.tex   (run twice, for TOC/LastPage)
%
\documentclass[11pt]{article}

\usepackage[T1]{fontenc}
\usepackage[utf8]{inputenc}
\usepackage[sc]{mathpazo}
\linespread{1.06}
\usepackage[letterpaper,margin=1in,headheight=15pt,headsep=18pt,footskip=32pt]{geometry}
\usepackage{microtype}
\usepackage{xcolor}
\usepackage{array}
\usepackage{booktabs}
\usepackage{tabularx}
\usepackage{longtable}
\usepackage{ragged2e}
\usepackage{enumitem}
\usepackage{fancyhdr}
\usepackage{lastpage}
\usepackage{hyperref}

% ------------------------------------------------------------------ colors --
\definecolor{byunavy}{HTML}{002E5D}   % BYU brand navy
\definecolor{byuroyal}{HTML}{0062B8}  % BYU royal (hyperlinks)
\definecolor{metagray}{gray}{0.42}

\hypersetup{
  colorlinks=true,
  linkcolor=byunavy,
  urlcolor=byuroyal,
  pdftitle={Getting Approval at BYU: The Steps and the COI Management Plan},
  pdfauthor={Vertical Cloud Lab (vertical-cloud-lab/byu-vcl)},
  pdfsubject={Integrated conflict-of-interest management plan and BYU approval sequence},
  pdfkeywords={BYU, conflict of interest, management plan, faculty founder, technology transfer}
}

% ------------------------------------------------------------------ layout --
\setlength{\parindent}{0pt}
\setlength{\parskip}{6pt plus 1pt minus 1pt}
\raggedbottom
\widowpenalty=9000
\clubpenalty=9000
\emergencystretch=1.5em
\setlist{itemsep=3pt,parsep=0pt,topsep=5pt,leftmargin=1.7em}

\pagestyle{fancy}
\fancyhf{}
\fancyhead[L]{\footnotesize\scshape\color{metagray} Conflict of Interest Management Plan --- Vertical Cloud Lab}
\fancyfoot[C]{\footnotesize\color{metagray} Page \thepage\ of \pageref{LastPage}}
\renewcommand{\headrulewidth}{0.3pt}
\renewcommand{\headrule}{{\color{metagray!50}\hrule width\headwidth height\headrulewidth}}
\fancypagestyle{plain}{%
  \fancyhf{}%
  \fancyfoot[C]{\footnotesize\color{metagray} Page \thepage\ of \pageref{LastPage}}%
  \renewcommand{\headrulewidth}{0pt}}

% ---------------------------------------------------------------- commands --
\makeatletter
% Top-level heading (Part 1, Part 2, closing section) with a navy rule.
\newcommand{\parthead}[1]{%
  \par\addvspace{22pt plus 6pt minus 4pt}%
  \phantomsection
  \addcontentsline{toc}{section}{#1}%
  \noindent{\color{byunavy}\Large\bfseries #1}%
  \par\nobreak\addvspace{5pt}%
  {\color{byunavy}\hrule height 1.1pt}%
  \par\nobreak\addvspace{10pt}\@afterheading}
% Second-level heading.
\newcommand{\subhead}[1]{%
  \par\addvspace{14pt plus 4pt minus 2pt}%
  \phantomsection
  \addcontentsline{toc}{subsection}{#1}%
  \noindent{\color{byunavy}\large\bfseries #1}%
  \par\nobreak\addvspace{6pt}\@afterheading}
% Numbered section of the management plan itself (renders as "S:N. Title").
\newcounter{psec}
\newcommand{\psection}[1]{%
  \par\addvspace{16pt plus 4pt minus 2pt}%
  \stepcounter{psec}%
  \phantomsection
  \addcontentsline{toc}{subsection}{\S\thepsec.\ #1}%
  \noindent{\color{byunavy}\large\bfseries \S\thepsec.\hspace{0.55em}#1}%
  \par\nobreak\addvspace{6pt}\@afterheading}
% A numbered clause ("3.1", "3.2", ...) with a hanging indent.
\newcommand{\clause}[1]{%
  \par\addvspace{5pt plus 1pt}%
  \noindent\hangindent=2.5em\hangafter=1
  {\color{byunavy}\bfseries #1}\hspace{0.7em}\ignorespaces}
\makeatother

% Placeholder marker, policy-source citation, signature line.
\newcommand{\ph}{{[}\,\textbullet\,{]}}
\newcommand{\policyref}[1]{{\small\itshape(#1)}}
\newcommand{\sigline}{\raisebox{-1pt}{\rule{1.5in}{0.4pt}}}

% Framed notice box.
\newcommand{\noticebox}[1]{%
  \par\medskip\noindent
  {\setlength{\fboxsep}{11pt}\setlength{\fboxrule}{0.8pt}%
   \fcolorbox{byunavy!45}{byunavy!4}{%
     \begin{minipage}{\dimexpr\textwidth-2\fboxsep-2\fboxrule\relax}
       \small #1
     \end{minipage}}}%
  \par\medskip}

\begin{document}

% ================================================================== title ==
\thispagestyle{plain}

{\color{byunavy}\hrule height 2.4pt}
\vspace{16pt}
{\huge\bfseries\color{byunavy} Getting Approval at BYU}\\[9pt]
{\LARGE The Steps and the COI Management Plan}
\vspace{10pt}

{\color{metagray}\small Vertical Cloud Lab \,\textperiodcentered\, 19 August 2026}
\vspace{6pt}

{\color{byunavy}\hrule height 0.8pt}
\vspace{14pt}

How a faculty-founder arrangement actually gets approved at BYU --- who signs
what, in what order --- and the central instrument, the
\textbf{conflict-of-interest management plan}, written for the VCL scenario
developed across \href{https://github.com/vertical-cloud-lab/byu-vcl/pull/156}{PR \#156}.
Companion to
\href{https://github.com/vertical-cloud-lab/byu-vcl/blob/claude/issue-155-20260708-0433/byu-vc-founder-rules.md}{BYU
rules on VC money and founder involvement}, which holds the underlying policy
analysis.

\noticebox{BYU publishes the operative disclosure instrument: the
\href{https://spo.byu.edu/0000018e-c41e-dbed-afaf-e6ffc57a0001/financial-disclosure-form-6-30-2023-docx}{\textbf{Financial
Disclosure Form}} and its
\href{https://spo.byu.edu/0000018e-c41e-dbed-afaf-e6ffc7110002/financial-disclosure-form-instructions-6-30-2023-docx}{\textbf{Instructions}}
(both 30 Jun 2023, on the \href{https://spo.byu.edu/policies}{RAO policies
page}). The form's Section~II asks whether ``a proposed conflict of interest
management plan is attached'' --- the plan is expected to arrive \emph{with}
the investigator's disclosure, and Part~2 below is written to serve as that
attachment. What BYU does not publish is a plan template; Part~2's structure
follows the management tools BYU's policies name, the TTO's process documents,
and the plan templates peer universities do publish (Iowa, Colorado State,
Nebraska).}

\renewcommand{\contentsname}{\normalsize\color{byunavy}Contents}
{\setlength{\parskip}{0pt}\small\tableofcontents}

% ================================================================= part 1 ==
\parthead{Part 1 --- The approval sequence}

Two clarifications up front, answering the questions as asked:

\begin{enumerate}[leftmargin=2em]
\item \textbf{``Department approval'' is the first layer of a stack, not a
single event.} What the department itself controls is steps 1--3 below: the
disclosure, the chair/dean written authorization of the outside role, and the
chair-written/dean-approved conflict management plan. Every transaction after
that (license, sponsored research agreement, board seat, leave) adds
university-level approvers on top --- but the chair and dean reappear as
signatories inside almost all of them, which is why the informal chair briefing
is step~0.
\item \textbf{Yes, a management plan is involved --- in fact BYU's policies
name two, and this scenario triggers both.} The
\href{https://policy.byu.edu/view/conflict-of-time-or-interest-policy}{Conflict
of Time or Interest Policy} (COTI) calls for a \textbf{conflict management
plan} when a disclosed outside activity creates a conflict or its appearance:
it is ``written with the chair, approved by the dean, and submitted through the
[Faculty Profile System],'' and may include ``any actions sufficient to avoid,
eliminate, or ensure against Conflicts or the appearance of Conflicts,
including recusal from discussions or decisions if a potential Conflict
exists.'' The
\href{https://policy.byu.edu/view/financial-conflict-of-interest-in-research-policy}{Financial
Conflict of Interest in Research Policy} (FCOI) calls for a \textbf{management
plan} when the AAVP--Research determines a disclosed Significant Financial
Interest is a conflict with respect to research; its named tools run from
public disclosure in publications and an independent monitor up to
disqualification from the project and ``reduction or elimination of the
financial interest (e.g., sale of an equity interest).'' Authorship and
approval of the FCOI plan are split by design: ``The Investigator works in
consultation with the RAO to develop the Management Plan. The AAVP--R has been
designated by the university to solicit, review, and approve Management
Plans.'' The COTI plan governs
the \emph{role and time} (co-founder title, days per month, resource wall); the
FCOI plan governs the \emph{research} (the company-sponsored work, the
students, the data). For a founder whose company licenses BYU IP and sponsors
research in their own lab, the sensible artifact is \textbf{one integrated plan
approved by chair, dean, and AAVP--R together} --- which is what Part~2
provides.
\end{enumerate}

\subhead{The steps in order}

{\footnotesize
\setlength{\extrarowheight}{2.5pt}
\begin{longtable}{@{} >{\bfseries}p{0.24in} >{\RaggedRight\arraybackslash}p{1.02in} >{\RaggedRight\arraybackslash}p{1.32in} >{\RaggedRight\arraybackslash}p{2.06in} >{\RaggedRight\arraybackslash}p{1.06in} @{}}
\toprule
{\color{byunavy}\scshape\#} &
{\color{byunavy}\scshape What happens} &
{\color{byunavy}\scshape Instrument} &
{\color{byunavy}\scshape Approver(s)} &
{\color{byunavy}\scshape When} \\
\midrule
\endfirsthead
\toprule
{\color{byunavy}\scshape\#} &
{\color{byunavy}\scshape What happens} &
{\color{byunavy}\scshape Instrument} &
{\color{byunavy}\scshape Approver(s)} &
{\color{byunavy}\scshape When} \\
\midrule
\endhead
\midrule
\multicolumn{5}{r@{}}{\itshape continued on next page}\\
\endfoot
\bottomrule
\endlastfoot
0 & Informal chair briefing & conversation + one-pager & --- & now, before the
company takes further shape \\
\addlinespace
1 & Disclose the outside activity & Faculty Outside Activities Form --- the
Conflict of Interest and Conflict of Time screen in
\href{https://facultyprofile.byu.edu/}{Faculty Profile} (Annual Review section)
& filed, not approved & \textbf{before engaging} in the activity; re-filed
annually \\
\addlinespace
2 & Authorize the role itself & prior \textbf{written} authorization &
department chair, dean, or supervisor (consulting: chair) & before taking the
co-founder/advisor role or signing anything company-side \\
\addlinespace
3 & COTI conflict management plan & plan \textbf{written with the chair} &
\textbf{dean approves}; submitted via Faculty Profile & as soon as a
conflict or its appearance is identified --- realistically, at step~2 \\
\addlinespace
4 & FCOI disclosure $\rightarrow$ determination $\rightarrow$ management plan
& \href{https://spo.byu.edu/policies}{Financial Disclosure Form}, with the
proposed management plan attached; plan developed by the investigator in
consultation with RAO & reviewed by chair $\rightarrow$ Dean's Office
$\rightarrow$ RAO director; \textbf{AAVP--Research} (Larry Howell) determines
FCOI status and approves all management plans & Form to the chair
\textbf{before proposal submission} to an external sponsor (internal funding:
before funds are expended); updates within 30 days of any new SFI; AAVP--R
review of mid-project SFIs within 60 days \\
\addlinespace
5a & License / assignment of BYU IP & option $\rightarrow$ term sheet
$\rightarrow$ license or assignment (TTO) & reviewed by chair, dean, TTO
director, General Counsel; \textbf{decision by AAVP--R}; founder recused from
negotiation & at deal time \\
\addlinespace
5b & Company-sponsored research in the lab & proposal + SRA (RAO) & chair +
dean approve the proposal; \textbf{AAVP--R signs} before it leaves BYU; RAO
negotiates opposite a non-founder company negotiator & when company money funds
lab work \\
\addlinespace
5c & Voting board seat at the licensee & written approval & dean + AAVP--R +
TTO --- and only if $\leq$\$3M raised, no non-accredited investors & rare; a
priced round forecloses it \\
\addlinespace
5d & Exception to standard income distribution (the $\geq$10\% forfeiture) &
written approval & dean + AAVP--R + TTO & if relief from the forfeiture/match
collision is sought \\
\addlinespace
5e & Personal leave (only if the posture changes) & leave application + written
leave agreement & chair $\rightarrow$ dean $\rightarrow$ AAVP--Faculty
Development & before any officer role \\
\addlinespace
5f & Company equipment placed in the lab & loan/bailment agreement or SRA term
& RAO signs (PI may not) & at placement \\
\addlinespace
5g & Selling instrument time to external users (cost recovery) &
recharge-center proposal + annual rate proposals
(\href{https://spo.byu.edu/00000190-3687-d777-a591-37af392d0001/rechargecenterpolicyprocedures-pdf}{Recharge
Center Policy and Procedures}) & chair $\rightarrow$ dean $\rightarrow$
Regulatory Accounting (rates approved annually) & before billing any external
user \\
\end{longtable}}

Notes on the table:

\begin{itemize}
\item \textbf{Step 1 timing is a policy command, not advice:} Outside
Activities ``must be disclosed \emph{before} engaging in them,'' then annually.
Since the April 2025 rewrite, the disclosure workflow lives inside the COTI
policy itself (Faculty Outside Activities Form; conflict management plans
through the same system) --- there is no separate HR procedure to hunt down.
\item \textbf{Step 2 is triggered by the role, not the money.} Written
authorization is required to ``operate a business, serve as an officer (in
title or in fact) of a business, or maintain employment with a business that
requires attention during the employee's assigned work hours,'' and Secondary
Employment (consulting) needs prior written chair approval within the
4-days/month cap. The test the approver applies is written down: chairs and
deans ``do not approve activities they believe would unduly limit a faculty
member's availability to students or compromise\ldots\ teaching, scholarship,
or citizenship performance'' --- which is why the one-pager for step~0 should
lead with the time budget and the teaching/mentoring plan, not the cap table.
\item \textbf{Step 4's instrument is public, and it expects the plan to come
from the investigator.} The \href{https://spo.byu.edu/policies}{Financial
Disclosure Form and its Instructions} (30 Jun 2023) are on the RAO policies
page. The form runs assessment (management position / SFI / sponsor-paid
travel) $\rightarrow$ disclosure (entity, nature, value bands from
``\$1--4,999'' up to ``$>$\$100K in \$50K increments'' or ``cannot be readily
determined'' --- the honest answer for private startup equity) $\rightarrow$
investigator certification $\rightarrow$ an \textbf{institutional signature
block of chair, dean/associate dean, RAO director, and AAVP--R}, plus an
RAO-use section tracking FCOI determination, interim measures, plan approval,
and sponsor reporting. Its Section~II line --- ``A proposed conflict of
interest management plan is attached: Yes~/~No'' --- is the hook Part~2 is
written for. The Instructions add the enforcement shape: once any \$5,000 SFI
threshold is met in a year, \emph{all} related SFIs are disclosed regardless
of value; new mid-project SFIs get an AAVP--R review within 60 days; and
untimely disclosure or plan noncompliance triggers a documented
\textbf{retrospective review} by the AAVP--R, chair, associate dean, and RAO
director. One wrinkle: the form's non-public-entity SFI definition uses the
\$5,000 aggregate test, while the policy also captures any equity interest;
the Instructions resolve it --- ``Any discrepancies between the Instructions
and the Policy will be resolved in favor of the Policy'' --- so founder equity
is disclosed regardless of value. Separately, all faculty on sponsored
research complete \href{https://spo.byu.edu/responsible-conduct-of-research}{CITI
COI training} with a refresher every 4 years --- a box the AAVP--R will expect
to be checked before approving anything.
\item \textbf{Steps 3 and 4 are where ``the management plan'' lives} --- see
the two-plans clarification above. For the license-only route Howell
recommended, it is possible the COTI layer resolves at step~2 with
authorization plus disclosure alone, and the FCOI plan only comes into
existence when the first proposal touching the company is submitted. But
Howell's own advice from the July 2026 meeting cuts the other way: get the
disclosures and the plan \textbf{dated early, while the stakes look small} ---
an approved plan is the artifact that answers later scrutiny; a retroactive one
is an admission.
\item \textbf{Step 5a includes the founder's chair by design.} The IP Policy:
``Any proposed transfer of Intellectual Property from the university will be
reviewed by the chair, dean, director of the Technology Transfer Office or
Creative Works Office, Office of General Counsel, and associate academic vice
president -- research and graduate studies,'' with the final decision the
AAVP--R's. The founder's sanctioned touchpoints during the deal are exactly
three: share information with TTO about the company as a potential licensee,
sign the license acknowledgment (``blue sheet''), and negotiate/sign the income
\textbf{Distribution Agreement} --- the one contract that is between BYU and
its own employee rather than with a third party, and therefore the founder's
only legitimate seat at any table.
\item \textbf{What to bring to step 0.} A one-page brief: what the company
does; proposed personal role and title (advisory, non-officer); equity range
and whether it will cross 10\%; what BYU touchpoints are contemplated
(license/assignment, BYU-held equity, later SRA, equipment placements); the
monthly time budget against the 4-day cap; how teaching, mentoring, and the
proposal pipeline are unaffected; and a draft management plan (Part~2) showing
the protections are already thought through. The chair's signature appears at
steps 2, 3, 5a, and 5b --- the person who should never be surprised.
\end{itemize}

\subhead{The affirmative case: how the arrangement enables the BYU portfolio}

The approval criterion chairs and deans apply is protective (availability to
students; teaching, scholarship, citizenship), so the step-0 brief and the
plan itself should also state what BYU \emph{gains}. For the VCL scenario,
concretely:

\begin{enumerate}[leftmargin=2em]
\item \textbf{Instrument access students could not otherwise have.}
Company-owned equipment co-hosted in the lab under an RAO-signed placement
(step~5f) gives research students hands-on access to instrumentation that
would be difficult to procure on startup-scale university budgets --- the same
bootstrap pattern as the CMU--Emerald Cloud Lab arrangement in the precedents
survey.
\item \textbf{Non-federal cost recovery that stays in the lab.} External users
billed for instrument time at the external recharge rate generate revenue with
far more spending flexibility than federal funds, because the
\href{https://spo.byu.edu/00000190-3687-d777-a591-37af392d0001/rechargecenterpolicyprocedures-pdf}{Recharge
Center Policy} allows ``surcharges\ldots\ to rates charged to non-University,
non-Federal agency users'' that ``can be used to reduce rates charged to
users, build a working capital reserve, or to finance equipment purchases,''
while wages of those operating the center are recoverable direct costs inside
any rate. Real constraints to plan around: external rates must include
institutional overhead unless Regulatory Accounting approves otherwise;
external revenues are tracked separately; above-cost external sales raise the
UBIT question the policy itself flags; the surcharge-use list does not name
student travel, so travel support should be papered through an allowable
channel (e.g., department funds freed by recharge recovery) with Regulatory
Accounting's sign-off. Self-imposed boundary, stated in the plan (\S5.6): none
of this revenue funds the faculty member's own salary or travel.
\item \textbf{Sponsored research and SBIR/STTR pathways.} Company sponsorship
converts to arm's-length SRAs (step~5b) with RAO negotiating --- new funding
into the lab, with SBIR/STTR subawards available later as a secondary channel
rather than a driving goal.
\item \textbf{Industry connections and student outcomes.} Downstream
application for the lab's research, industry contacts for the group, and
post-graduation employment opportunities for students at the company and its
partners.
\item \textbf{Open source keeps student publication paths unobstructed.}
Because the company's model is open software and open hardware, student work
touching it flows to manuscripts and conference talks with only the TTO
disclosure gate in the way --- where sponsorship by a conventional proprietary
company typically adds confidentiality reviews, publication delays, and thesis
complications.
\end{enumerate}

% ================================================================= part 2 ==
\parthead{Part 2 --- The integrated conflict-of-interest management plan}

The plan below is written for the specific posture the PR \#156 thread
converged on after the Howell meeting: \textbf{stay full-time faculty (no
leave), listed co-founder with equity, advisory non-officer title, TTO license
or assignment with BYU possibly taking equity, the company later sponsoring
research in the founder's own lab, open-source-by-default outputs, possible
company equipment in the lab, and students throughout.} Bracketed text is
placeholder; provisions cite their policy source so approvers can check
provenance. Provisions marked \emph{[snapshot]} rest on login-gated policies
reviewed in the
\href{https://github.com/vertical-cloud-lab/byu-vcl/pull/156\#issuecomment-5094216608}{July
2026 full-policy-snapshot cross-check} rather than on publicly fetchable text.

One structural point worth understanding before reading it: the snapshot
cross-check found that \textbf{BYU has no default student-protection rule for
faculty startups} (no analogue of Stanford's flat bar on involving current
students in company activities). At BYU those protections exist \emph{only if a
management plan creates them} --- which is why \S6 below is the longest
section.

% ------------------------------------------------- the plan document itself --
\clearpage
\vspace{6pt}

\begin{center}
{\LARGE\bfseries\color{byunavy} Conflict of Interest Management Plan}\\[7pt]
{\normalsize\bfseries Brigham Young University --- Department of Mechanical
Engineering / Ira A.\ Fulton College of Engineering}\\[10pt]
{\color{byunavy}\hrule height 0.8pt}
\end{center}
\vspace{4pt}

{\setlength{\extrarowheight}{3pt}
\begin{tabular}{@{} >{\leavevmode\color{byunavy}\scshape}p{1.45in} >{\RaggedRight\arraybackslash}p{4.55in} @{}}
Faculty member & [Sterling G. Baird], Assistant Professor, Department of
Mechanical Engineering (the ``Faculty Member'') \\
Outside entity & [NewCo, Inc.], a [Delaware C-corporation] commercializing
automated (``cloud lab'') materials-experimentation technology (the
``Company'') \\
Governing policies & Conflict of Time or Interest Policy (7 Apr 2025);
Financial Conflict of Interest in Research Policy (7 Apr 2025); Financial
Disclosure Form and Instructions (30 Jun 2023); Intellectual
Property Policy (rev.\ 16 Dec 2024) and Intellectual Property Procedures;
Sponsored Projects Handbook; Recharge Center Policy and Procedures \\
Plan effective date & [date] \\
Review cycle & annual, with the Faculty Outside Activities Form \\
\end{tabular}}

\psection{Disclosed facts and interests}

Per the Faculty Outside Activities Form dated \ph\ and the FCOI disclosure
dated \ph:

\clause{1.1} The Faculty Member is a \textbf{co-founder} of the Company.

\clause{1.2} The Faculty Member holds an equity interest of \textbf{\ph\%
fully diluted} ([below 10\% / exactly 10\% / above 10\%] --- the provisions of
\S9 apply accordingly). No salary, consulting fees, or other remuneration are
received from the Company (see \S3.4).

\clause{1.3} The Faculty Member's role is \textbf{[co-founder and Chief
Scientific Advisor]} --- an advisory role with no line-management authority, no
employment relationship, and no officer title, in fact or in name.

\clause{1.4} The Company [is negotiating / holds] a [license / assignment] to
BYU-owned intellectual property arising from the Faculty Member's laboratory
(TTO ref.\ \ph). BYU [holds / may accept] equity in the Company as partial
consideration.

\clause{1.5} The Company [intends to / may] sponsor research in the Faculty
Member's BYU laboratory under a sponsored research agreement (``SRA'')
negotiated by the Research Administration Office (``RAO'').

\clause{1.6} The Company may place Company-owned equipment in the laboratory
under a university-signed agreement (\S7.2).

\clause{1.7} Students and postdoctoral researchers work in the laboratory;
some may be funded under the SRA.

\clause{1.8} The laboratory's stated policy is open-source release of outputs
by default, subject to the TTO disclosure gate (\S7.3).

\clause{1.9} The Faculty Member remains on a full-time faculty appointment and
is not on leave. A change in any fact in this section is a plan-amendment
trigger (\S10).

\psection{Determinations}

\clause{2.1} The department chair and dean have authorized the outside
activity described in \S1.1--1.3 in writing, subject to this plan.
\policyref{COTI: written authorization to ``operate a business, serve as an
officer (in title or in fact) of a business\ldots''}

\clause{2.2} The AAVP--Research has determined that the interests in \S1
constitute a Financial Conflict of Interest with respect to [the SRA in \S1.5
and any research related to the licensed IP], and that the conflict can be
\textbf{managed} under this plan rather than requiring elimination.
\policyref{FCOI: AAVP--R determination; management-plan tools.}

\clause{2.3} This plan is intended to manage conflicts \textbf{and the
appearance of conflicts}. \policyref{COTI management-plan standard.}

\clause{2.4} The parties record the university benefits anticipated from the
arrangement, alongside the conflicts it manages: research-student access to
Company-hosted instrumentation the university would find difficult to procure
(\S7.2); non-federal cost-recovery and sponsored-research revenue to the
laboratory (\S5.1, \S5.6); industry connections, downstream application of
the laboratory's research, and SBIR/STTR pathways; post-graduation employment
opportunities for students; and open-source outputs that keep student
publication and conference paths unobstructed (\S1.8, \S7.3--7.4). These
benefits motivate the approvals; they do not relax the protections below.

\psection{Role, title, and time}

\clause{3.1} The Faculty Member will hold \textbf{advisory titles only}
(e.g., ``Chief Scientific Advisor,'' scientific advisory board chair) and will
not serve as an officer of the Company, in title or in fact, nor as a line
manager or employee, while a full-time faculty member not on leave.
\policyref{COTI officer clause; Sponsored Projects Handbook: ``Faculty members
not on leave are not to be line officers in outside companies.''}

\clause{3.2} The Faculty Member will not hold a \textbf{voting board seat}.
If the Company has raised $\leq$\$3{,}000{,}000 with no non-accredited
investors and a voting seat is desired, it requires separate prior written
approval of the dean, AAVP--R, and TTO; the parties acknowledge a priced
venture round forecloses it, and a board-observer or advisory seat is the
durable shape. \policyref{IP Policy board-seat rule.}

\clause{3.3} Aggregate time on Company matters will adhere strictly to the
\textbf{four days per month} maximum while under contract, averaged as one day
per week across all outside activities combined. Any Company time spent during
the workday is made up with equivalent time on adjacent evenings and weekends.
Hours are logged and reported with the annual disclosure. \policyref{COTI
Secondary Employment cap and compensating-time expectation; Expectations of a
Faculty Appointment aggregation.}

\clause{3.4} \textbf{No dual compensation channel.} While the Faculty Member
participates in Company-sponsored research at BYU (\S5), the Faculty Member
will receive \textbf{no consulting fees or other paid-services compensation
from the Company} --- founder equity is the sole financial interest, managed
under this plan. Any future paid consulting arrangement requires prior written
chair approval and either (a) withdrawal from participation in
Company-sponsored BYU research, or (b) a written AAVP--R exception under the
rare-exception clause. \policyref{IP Policy: a faculty member ``may either act
as a paid consultant to or participate in sponsored research for the same
company but may not perform both at the same time,'' rare exceptions by
approved management plan + written AAVP--R approval.}

\clause{3.5} \textbf{Teaching and mentoring take scheduling priority.} Company
matters are never scheduled against, and never displace, scheduled class time,
class preparation, student meetings, or research meetings; scheduled teaching
obligations are categorically outside the \S3.3 allowance. In calendaring,
meetings with students and research meetings are placed first and Company
commitments fit around them. \policyref{COTI approval criterion: activities
may not ``unduly limit a faculty member's availability to students or
compromise\ldots\ teaching, scholarship, or citizenship performance.''}

\psection{Contracting wall (recusal)}

\clause{4.1} The Faculty Member will not \textbf{negotiate, influence, or
attempt to influence} the negotiation of any contract or agreement between BYU
and the Company --- including the license/assignment, any SRA, any equipment
loan, and any gift --- \textbf{from either side of the table}. \policyref{COTI
Related Party clause, applied as a plan term regardless of the \S9 threshold.}

\clause{4.2} Company-side negotiations with BYU are conducted by [non-founder
officer/business lead, name/title]; BYU-side negotiations only by RAO, TTO, or
the Office of General Counsel per their signature authority.

\clause{4.3} The Faculty Member will not advocate acceptance or terms of any
BYU--Company agreement with the chair, dean, RAO, TTO, AAVP--R, or through
colleagues.

\clause{4.4} \textbf{Permitted activities} (the sanctioned channels,
exhaustively): (a) sharing information with TTO about the Company as a
potential licensee; (b) technical input to the SRA statement of work, submitted
through RAO's designated contact; (c) serving as the Company's named technical
contact after execution of an agreement; (d) reviewing and signing the license
acknowledgment (``blue sheet''); (e) negotiating and signing the income
\textbf{Distribution Agreement} with BYU, which is an internal BYU--employee
contract outside the Related Party prohibition. \policyref{TTO ``Working with
the Technology Transfer Office''; COTI.}

\psection{Company-sponsored research in the laboratory}

\clause{5.1} Any Company use of laboratory personnel, students, or facilities
is converted to an \textbf{SRA with cost reimbursement} at BYU's standard terms
(51.5\% MTDC indirect on the MTDC base; no permanent publication restrictions;
TTO-reviewed IP terms). No Company work product, data generation, or product
development occurs outside the SRA's statement of work. \policyref{COTI
resource wall; Sponsored Projects Handbook.}

\clause{5.2} \textbf{Project leadership:} [Option A --- Prof.\ \ph, who holds
no Company interest, serves as PI of record; the Faculty Member participates as
co-investigator.] [Option B --- the Faculty Member serves as PI, with the
independent oversight of \S5.3 and enhanced monitoring of \S8.] If the Faculty
Member ever takes leave, a genuinely independent interim PI holding a regular
faculty appointment is mandatory (research associates and postdocs may not
serve as PI \emph{[snapshot: Non-CFS Track Academic Appointments Policy]}), and
the Faculty Member's residual role is limited to a disclosed advisory capacity
named in the leave agreement.

\clause{5.3} \textbf{Independent oversight of the science:} Prof.\ \ph\ (no
Company interest) reviews the design, data integrity, and conclusions of
Company-sponsored work [annually / prior to submission of each publication];
disagreements are escalated to the chair. \policyref{FCOI management tools:
independent monitor, modification of research plan.}

\clause{5.4} Publications and presentations arising from Company-sponsored or
license-related work disclose the Faculty Member's Company interest and, if
applicable, BYU's equity position. \policyref{FCOI public-disclosure tool;
institutional-COI practice.}

\clause{5.5} Student publication rights and the no-permanent-restriction
baseline are preserved in every SRA. Sponsor prepublication review is limited
to short delays for proprietary-information and patent screening.

\clause{5.6} \textbf{Cost recovery for instrument time, and where that money
goes.} Any sale of laboratory instrument time or services to the Company or
other external users outside an SRA occurs only through a
recharge/service-center account established under the Recharge Center Policy
and Procedures (chair and dean approve the proposal and rates; Regulatory
Accounting approves accounts and annual rate proposals). External users are
billed at the external rate --- institutional overhead included, plus any
approved surcharge --- with external revenues tracked separately; surcharge
proceeds are applied as the policy permits (reducing rates, working-capital
reserve, financing equipment purchases), and the UBIT question on above-cost
external sales is resolved with Regulatory Accounting before billing begins.
\textbf{Self-imposed boundary, beyond what policy requires:} revenue
attributable to Company-related cost recovery or sponsorship is budgeted to
student wages, equipment, and student travel or other laboratory purposes ---
none of it funds the Faculty Member's salary (including summer salary) or the
Faculty Member's own travel. \policyref{Recharge Center Policy and Procedures;
surcharge-use list quoted therein.}

\psection{Students and postdoctoral researchers}

\clause{6.1} \textbf{Written notice.} Every member of the laboratory receives
written notice of the Faculty Member's Company roles and interests (and BYU's,
if BYU holds equity) at onboarding and annually; the notice is copied to the
chair. Students assigned to Company-sponsored projects acknowledge it in
writing.

\clause{6.2} \textbf{Independent advocate.} The chair (or a designee holding
no Company interest) meets each student working on Company-sponsored or
license-related projects at least once per [semester], without the Faculty
Member present, to review degree progress, topic independence, and any
concerns; students may also raise concerns directly with the graduate
coordinator or dean's office at any time.

\clause{6.3} \textbf{No company work outside the SRA.} No student is directed
to perform services for the Company except within the SRA's statement of work.
Using students ``to perform services for an outside entity'' in which the
investigator holds an interest is a listed FCOI conflict; this plan resolves it
by containment within the SRA plus the protections of this section.

\clause{6.4} \textbf{Company employment of students.} A student who wishes to
intern with or join the Company is referred to the [chair/dean's office] for
independent advice before any offer is accepted (the Stanford-model referral).
The Faculty Member will not simultaneously supervise the same student at BYU
and at the Company; if an exception is unavoidable, the student's thesis
committee must carry a majority of members with no Company interest and the
dual role is disclosed to the committee in writing. Combined BYU + Company
employment respects student hour caps (20 hours/week during semesters), with
particular attention to international students' visa limits. \emph{[snapshot:
Student Hiring and Employment Policy]}

\clause{6.5} \textbf{Theses and timing.} Thesis/dissertation scope and defense
timing are set by the student and committee on academic grounds; no embargo
(ADV Form 8e) or defense NDA is requested for the Company's commercial
convenience --- where patent timing requires confidentiality, the default is a
provisional filing before the defense, not securing the thesis.

\clause{6.6} \textbf{IP and credit.} Lab participants sign the Employee and
Student Assignment of Ownership and Non-Disclosure Agreement (TT-1089) at
onboarding. Students are informed in writing that inventorship is a legal
determination based on conception (made during prosecution, not by the lab),
that ``Developer'' royalty status and Company compensation are separate
questions, and how the lab recognizes contribution (authorship; Distribution
Agreement shares where inventorship supports it).

\psection{University resources, equipment, and intellectual property}

\clause{7.1} \textbf{Resource wall, with a defined office allowance and
complete digital separation.} No university resources, facilities, work time,
personnel, or IT accounts are used for Company operations. All Company work
runs on Company-owned laptops, accounts, software subscriptions, and storage
from day one --- complete separation from BYU devices, site licenses (licensed
for university education/research use, not commercial use), and IT accounts,
which additionally require responsible-VP + CIO written authorization for any
non-university commercial use \emph{[snapshot: Appropriate Use of IT Resources
Policy]}. The chair's written authorization under this plan extends to
\textbf{incidental use of the campus office} for Company communications
(calls, correspondence, scheduling) conducted on Company devices, with the
time made up under \S3.3; Company operations, investor or customer meetings,
records storage, and any use of university personnel remain off campus and
outside BYU systems. \policyref{COTI: ``University Personnel may not use
university resources, facilities, or work time to engage in commercial
activities, businesses, or `start-up' companies''; Expectations of a Faculty
Appointment requires written approval to use a campus office for business
purposes [snapshot] --- this clause is that written approval, scoped.}

\clause{7.2} \textbf{Equipment placements.} Company-owned equipment enters the
laboratory only under a university-signed instrument --- an RAO equipment
loan/bailment agreement or an SRA term --- settling title, insurance,
maintenance, access, export-control screening, and permitted use before
delivery. The Faculty Member signs nothing. Students operate such equipment
only for BYU research or within the SRA. In-kind resources are reported on
federal awards (Facilities \& Other Resources / Other Support) as required.

\clause{7.3} \textbf{TTO gate before public release.} All potentially
protectable laboratory outputs are disclosed to TTO \textbf{before any public
disclosure} --- public repositories, preprints, demos, and public theses
included --- and released after TTO either releases the rights or files a
provisional. The laboratory maintains a standing pre-release checklist, and the
parties will agree on a target turnaround with TTO so the open-source release
cadence is plannable. \policyref{IP Policy disclosure-before-public-use; TTO
``Preparing a Patent Application.''}

\clause{7.4} \textbf{Open-source boundaries are set upstream, not at the
table.} The laboratory's open-science policy is public, dated, and
sponsor-agnostic; the Company's open/proprietary split lives in the Company's
published licensing policy, adopted by the Company outside any BYU negotiation.
The Faculty Member does not advocate particular IP or release terms in any
BYU--Company negotiation (\S4).

\clause{7.5} \textbf{Invention hygiene and the prong-(c) acknowledgment.}
Records are maintained sufficient to distinguish BYU-side from Company-side
conception; joint-inventorship questions go to TTO. The Faculty Member
acknowledges the IP Policy's ownership claim over IP ``related to the current
or demonstrably anticipated business, research, or development of the
university'' and will obtain \textbf{written TTO releases or license
carve-outs} for Company-side inventions rather than relying on
time-and-resource records alone.

\psection{Transparency, certification, and monitoring}

\clause{8.1} The Company relationship is disclosed in proposals, in
publications and presentations arising from affected research, to affected
students (\S6.1), and to journals/conferences per venue policy.

\clause{8.2} Annual re-disclosure via the Faculty Outside Activities Form;
FCOI disclosure updates within \textbf{30 days} of any new or changed
significant financial interest, including SFIs independent of existing
agreements or proposals. Once any \$5,000 SFI threshold is met in a year, all
further related SFIs are disclosed regardless of value. \policyref{Financial
Disclosure Form Instructions.}

\clause{8.3} The Faculty Member notifies the chair and AAVP--R within 30 days
of: any Company financing event; any change in title or governance role; equity
crossing the 10\% line in either direction (\S9.3); any new proposed
BYU--Company agreement; any contemplated leave.

\clause{8.4} An annual one-page compliance report under this plan goes to the
chair, dean, and AAVP--R with the annual disclosure; the AAVP--R may direct an
independent monitor at any time. \policyref{FCOI tools.}

\clause{8.5} On PHS-funded research, the Faculty Member acknowledges BYU will
respond within 5 business days to public requests for information about
disclosed SFIs --- name, title, entity, nature, and approximate value.

\psection{Equity-threshold provisions}

\clause{9.1} If the Faculty Member's interest is \textbf{ten percent or more},
the Developer distribution under the IP Policy (the 45\% share, and with it the
research-account designation and college/IPS match) is forfeited by default;
any relief requires a written \textbf{Exception to Standard Income
Distribution} approved by the dean, AAVP--R, and TTO. The Faculty Member may
request such an exception through the chair but will not lobby for it
(\S4.3).

\clause{9.2} If the interest \textbf{exceeds ten percent}, the Company is a
``Related Party'' under COTI and the \S4 recusal is mandatory policy, not
merely a plan term.

\clause{9.3} The Faculty Member provides the AAVP--R an updated capitalization
summary after each financing event, so threshold status is always current.

\psection{Change-of-posture triggers (amend the plan \emph{before} the change)}

Prior plan amendment and re-approval are required before:

\begin{itemize}
\item taking any officer title or board seat;
\item taking any leave;
\item becoming PI [or ceasing to use an independent PI] on Company-sponsored
work;
\item accepting any Company compensation beyond equity;
\item the Company hiring any student currently supervised by the Faculty
Member;
\item BYU units becoming customers of the Company (which separately triggers
the Employee-Vendor Policy's pre-approval machinery \emph{[snapshot]});
\item any human-subjects research touching the Company relationship.
\end{itemize}

\psection{Duration, amendment, and remedies}

\clause{11.1} This plan is reviewed annually by the AAVP--R with the chair and
dean, and may be amended, strengthened (independent monitor, disqualification,
divestiture), or terminated by the AAVP--R as circumstances change.
\policyref{FCOI tools.}

\clause{11.2} The plan terminates when the relationships and interests in \S1
end or fall below disclosure thresholds, confirmed in writing by the AAVP--R.

\clause{11.3} Noncompliance is handled under the governing policies, including
retrospective review and reporting obligations on federally funded work.

\psection{Acknowledgment and approvals}

{\setlength{\extrarowheight}{3pt}
\begin{tabularx}{\textwidth}{@{} >{\RaggedRight\arraybackslash}p{2.45in} >{\RaggedRight\arraybackslash}X >{\centering\arraybackslash}p{1.62in} @{}}
\toprule
{\color{byunavy}\scshape Party} &
{\color{byunavy}\scshape Role under this plan} &
{\color{byunavy}\scshape Signature / date} \\
\midrule
{[Sterling G. Baird]}, Faculty Member & subject of the plan & \sigline \\[1.6em]
\ph, Department Chair, Mechanical Engineering & plan co-author (COTI); reviewer
of IP transfers & \sigline \\[1.6em]
\ph, Dean (or Associate Dean), Ira A.\ Fulton College of Engineering & plan
approval (COTI); Dean's Office disclosure review; board-seat/income-exception
approvals & \sigline \\[1.6em]
{[Debbie Silversmith]}, Director, Research Administration Office & disclosure
review; plan development with the Investigator (FCOI); negotiates any
SRA/equipment agreement (\S\S4--5, 7.2) & \sigline \\[1.6em]
{[Larry Howell]}, AAVP--Research \& Graduate Studies & FCOI determination and
plan approval; final IP-transfer decision & \sigline \\[1.6em]
{[Dave Brown]}, Director, Technology Transfer Office & for provisions touching
the license/assignment (\S\S4, 7, 9) & \sigline \\[1.6em]
\bottomrule
\end{tabularx}}

% ================================================================ closing ==
\clearpage
\parthead{Verification notes}

\begin{itemize}
\item \textbf{How this aligns with what BYU publishes.} The disclosure
workflow above is sourced from BYU's own operative instrument --- the
Financial Disclosure Form and Instructions on the
\href{https://spo.byu.edu/policies}{RAO policies page} --- and Part~2 is
structured to be the ``proposed conflict of interest management plan'' the
form asks the investigator to attach. What remains internal to BYU: any plan
boilerplate the AAVP--R office keeps, and the Conflict of Interest and
Conflict of Time screen inside
\href{https://facultyprofile.byu.edu/}{Faculty Profile} (login; help site at
\href{https://facultyprofilehelp.byu.edu/}{facultyprofilehelp.byu.edu}).
Asking the AAVP--R office whether a preferred format exists --- and how the
Warnick/Linear Signal founder-sponsor case was papered --- remains the
cheapest calibration step.
\item \textbf{Some provisions cite gated policies} reviewed only in the July
2026 policy-snapshot cross-check (marked \emph{[snapshot]}): the Non-CFS PI
restriction, student hour caps, the IT commercial-use authorization, the
campus-office clause, and the Employee-Vendor Policy. Verify current text with
a BYU login before relying on the details.
\item \textbf{The plan deliberately over-includes.} Provisions like \S5.3
(independent science oversight), \S5.6's self-imposed budget boundary, and
parts of \S6 exceed what BYU's written policies strictly require --- they
import the elements peer plans use and the FCOI policy's optional tools. That
is intentional (the plan is the artifact that answers scrutiny when things go
wildly well), but each one is negotiable with the approvers.
\end{itemize}

\vfill
{\color{metagray}\footnotesize\noindent
Typeset 19 August 2026 from
\href{https://github.com/vertical-cloud-lab/byu-vcl/blob/claude/issue-155-20260708-0433/byu-coi-management-plan-example.md}{\texttt{byu-coi-management-plan-example.md}}
(vertical-cloud-lab/byu-vcl, PR \#156). LaTeX source:
\texttt{latex/byu-coi-management-plan-example.tex} in the same repository.}

\end{document}
